\chapter*{How to Read This Book}
\label{ch:map}
\addcontentsline{toc}{chapter}{How to Read This Book}
\markboth{How to Read This Book}{How to Read This Book}
\renewcommand{\thefigure}{\Alph{figure}}\setcounter{figure}{0}
% ===================================================================

This book is long, and it is long in an uneven way: some of it is
confessional, some of it is a mathematics reference, and some of it
is a story about whales.
You are entitled to know the shape of the thing before you commit,
and to know which parts you can skip without losing the argument.
That is what this chapter is for.
The confessional part of the Pledge has just ended; this chapter and
the one after it are the only chapters in the book that are about the
book.

\section*{What it is trying to do}

One question runs underneath all of it.

\begin{quote}
What is the minimum mathematical structure needed to ground
physics, knowledge, life, and mind?
\end{quote}

\noindent
The answer developed here is that a \emph{graph-structured lambda
theory} --- a grammar of terms, some equations on those terms, and
some rules for rewriting them --- is already enough, provided you
add two things: a site at which two terms meet, and a meter on that
site.
Everything else in the Turn is a fact about a theory so equipped.

The three acts do different jobs.
The Pledge, which you are nearly finished with, is the ordinary
world: what computation is, why the version of it you are carrying
around is probably a function machine and needs updating, and why
the question ``what is it like to be a computer program'' is a
better question than the one usually asked.
The Turn is the mathematics, and it is most of the book.
The Prestige brings it back --- and it brings it back by explaining
how the trick was done.

i want to say that last part plainly here, at the front, because it
changes how the middle should be read.
The Prestige opens by showing you the method: the whole
construction was bought with a single restriction of scope, the
restriction yielded an ontology and a grounded language for free,
and Chapter~\ref{ch:trick} names the exact line where it happens.
The book audits itself before it finishes.
If you would rather know that going in than be surprised by it,
now you do.

\section*{The shape of it}

The Turn is divided into six parts of its own.
They do not depend on each other in the order they are printed,
and Figure~\ref{fig:map} records what actually depends on what.
The arrows are dependencies, not reading order.

\begin{figure}[htbp]
\centering
\begin{tikzpicture}[>=Stealth, node distance=1cm,
  box/.style={draw, rounded corners, minimum height=0.8cm,
              minimum width=3.5cm, align=center, font=\small},
  dep/.style={->, thick, gray!70}]

\node[box, fill=remarkgray] (mach) at (0,4.2) {The Machinery};
\node[box, fill=remarkgray] (eco)  at (0,2.8) {Ecology as\\Cognitive Architecture};
\node[box] (phys) at (-4.4,1.2) {Situating\\the Learner};
\node[box] (sci)  at (0,1.2) {Why Scientists\\Get Misled};
\node[box] (ent)  at (-4.4,-0.4) {Suffering,\\Homeostasis, Dissipation};
\node[box] (orig) at (0,-0.4) {Origins of Life};
\node[box, fill=remarkgray] (pre) at (0,-2.0) {The Prestige};
\node[box, fill=remarkgray] (pre2) at (0,-3.4) {\emph{i} the break from Turing\\\emph{ii} reference and metaphysics};
\draw[dep] (pre) -- (pre2);

\draw[dep] (mach) -- (eco);
\draw[dep] (mach.west) to[out=180,in=90] (phys.north);
\draw[dep] (eco) -- (sci);
\draw[dep] (phys) -- (ent);
\draw[dep] (eco.west) to[out=200,in=20] (ent.north east);
\draw[dep] (sci) -- (orig);
\draw[dep] (ent) -- (orig);
\draw[dep] (orig) -- (pre);
\draw[dep] (eco.east) to[out=340,in=20] (pre.east);

\node[font=\small\itshape, align=left, text width=4.6cm]
  at (5.3,2.4) {Shaded: the spine.\\
  Everything else can be\\ skipped on a first pass\\ without losing the\\ argument.};

\node[font=\small\itshape, align=left, text width=4.6cm]
  at (5.3,-0.6) {The two ladders are\\ perpendicular: one climbs\\
  by adding axioms, the\\ other by adding oracles.\\
  Neither is above the other.};

\end{tikzpicture}
\caption{What depends on what. The Machinery is used everywhere and
is a reference rather than a narrative. Ecology is the argumentative
core. The Prestige needs only those two.}
\label{fig:map}
\end{figure}

Two things about the figure are worth saying in words.

The first is that the shaded boxes are the spine.
If you read the Pledge, the Machinery, the Ecology, and the
Prestige, you have read the argument.
The other four parts of the Turn are the argument applied --- to
physics, to the limits of science, to what a learner gains by watching
its own reserves, to the origin of life --- and each is worth having, but
none is load-bearing for the others' conclusions.
The Prestige is not in that category: since the last revision it carries
the break from Turing and the account of consciousness as well as the
reveal, and it is now about a third again as long as it was.

The second concerns the two ladders, which are the single most
confusing structural feature of the book and which i have watched
more than one reader run aground on.
\emph{Situating the Learner} climbs a ladder of subcategories: each
chapter imposes one more axiom on a cost-accounted theory, and each
axiom carves out the theories that satisfy it.
The Prestige's first half climbs a different ladder, of choice
strength.
These are perpendicular.
A theory can be high on one and low on the other, and neither
ladder is a refinement of the other.
Remark~\ref{rmk:two-ladders} says this once, in passing, three
hundred pages in.
It should have been said here, so here it is.

\paragraph{A confusing detail about the part numbers.}
\LaTeX{} numbers every \texttt{part} in one sequence, so the Turn's
six internal parts print as Parts~III through~VIII, while the prose
--- following the Overarching Introduction --- calls them Parts~I
through~VI of the Turn.
The Prestige now has two internal parts of its own, which print as
Parts~X and~XI and which the prose calls its first and second halves.
Both conventions are used.
Where it matters i have referred to parts by name rather than by
number, and i recommend you do the same.

\section*{Five routes through}

\paragraph{Front to back.}
About six hundred pages.
This is the order the argument was built in and the order in which
each part earns the next, and if you have the time it is what i
would want.
Everything below is a compromise.

\paragraph{The short route --- about ninety pages.}
The Pledge, then Chapter~\ref{ch:sci} (\emph{The Mortal
Scientist}), then the Prestige.
That is the argument end to end: what a learner is, what it costs
to know anything, and what the book concludes about the whole
construction.
\emph{What you lose:} every reason to believe the machinery is well
founded, and the entire empirical contact surface.
You will have to take a good deal on trust.

\paragraph{The philosopher's route --- about two hundred pages.}
The Pledge, the Ecology part entire (Chapters~\ref{ch:sci}
to~\ref{ch:eng}), \emph{Why Scientists Get Misled}, then the
Prestige.
This is the epistemology without the physics: what a budgeted
knower can find out, why populations of knowers converge on a false
ontology, and what that does to the theory of knowledge.
\emph{What you lose:} the derivation of causal structure and
conservation.
You keep the argument that a mind can do something a Turing machine
cannot, since that argument now lives in the Prestige.

\paragraph{The mathematician's route.}
Skip the Pledge --- it is confessional and it is my burden, not
yours --- and start at the Machinery.
Then Ecology, then \emph{Situating the Learner}, then the Prestige,
whose first half is where the break from Turing was moved to and is the
part of this book a mathematician is most likely to want.
\emph{What you lose:} the motivation for every choice made in the
Machinery, which is all in the Pledge.
That is more of a loss than it sounds.

\paragraph{The origins route.}
Chapters~\ref{ch:sci} and~\ref{ch:eng}, then the Origins of Life
part entire.
For readers who came for assembly theory and want the argument that
bears on it.
\emph{What you lose:} why an ecology of learners is the right model
of a mind in the first place, which is the premise the whole
assembly-index argument runs on.

\section*{How hard it gets, and where}

The difficulty is not monotone and it is fair to say so.

The Pledge assumes nothing.
The Prestige is now two parts and they differ sharply.
Its first half --- the break from Turing, apertures, agency, choice
types, consciousness --- is genuinely technical and assumes the
Machinery.
Its second half assumes nothing until Chapter~\ref{ch:notation}, which
is the most technical thing there and which can be read for its
conclusions alone; Chapter~\ref{ch:sim} after it needs only
Chapter~\ref{ch:notation}'s two conditions and the idea of a lattice
position.

The Machinery is the driest part of the book, deliberately.
It is a \emph{reference}: it develops constructions once so that
everything after can use them without rebuilding them.
It is entirely legitimate to skim it, start the Ecology part, and
come back when a construction is used and you want to know what it
is.
That is how i would read it if i had not written it.

The Ecology part is the most worked and, i think, the most
rewarding.
It is also the only place in the book where a claim is checked
against a number.

\emph{Situating the Learner} and the Prestige's first half are the most
speculative, and both say so in their own closing chapters ---
Chapter~\ref{sec:hyper-gaps} at some length, since that half rests on
the least-built construction in the book.
\emph{Origins of Life} has the most contact with empirical
science and correspondingly the most exposure.

As for background: nothing requires category theory until
Chapter~\ref{ch:catgslt}, and that chapter builds what it needs.
Process-calculus vocabulary is introduced in
Chapter~\ref{ch:gslt} and used throughout.
Modal logic appears in Chapter~\ref{sec:lts-hml} and is developed
from scratch.
No physics background is assumed anywhere, and where physics is
invoked it is invoked as a target to be recovered rather than as a
premise.

\section*{How to read a claim}

The book states things, in its own phrase, precisely enough to be
wrong.
The typographic conventions are how you tell which kind of thing
you are looking at, and they are load-bearing.

\begin{itemize}[leftmargin=1.6em, itemsep=3pt]
\item \textbf{Theorem}, \textbf{Proposition}, \textbf{Lemma},
      \textbf{Corollary} --- proved here or in a cited source.
\item \textbf{Conjecture} --- believed, not proved, and named so
      that someone can settle it.
\item \textbf{Observation} --- true and usually easy, but worth
      having a number so it can be pointed at.
\item \textbf{Question} --- open, and stated in a form where an
      answer would be recognisable.
\item \textbf{Remark} --- commentary, including several remarks
      whose job is to say that something earlier was wrong.
\end{itemize}

\noindent
In the two chapters that draw correspondences with physics, the
tables carry a status column reading \textsc{con}, \textsc{thm}, or
\textsc{cnj} --- construction, theorem, conjecture --- because a
table of correspondences is exactly the place where a reader can
lose track of which rows are earned.

\subsection*{Where the weight rests on something unproved}

This is the part of the chapter i would most want a sceptical
reader to have.
The following claims are load-bearing and not settled.
They are flagged where they appear, but they appear scattered
across six hundred pages, and a reader who discovers them one at a
time will reasonably wonder what else is hiding.
Nothing else is hiding.

\begin{itemize}[leftmargin=1.6em, itemsep=4pt]
\item Conjecture~\ref{thm:compact-consensus}, on consensus under
      compactness, was stated as a theorem in earlier drafts and the
      proof does not work.
      It is now a conjecture, with the failed argument and both
      failed repairs written out beside it, because the way it fails
      says something: compactness is a richness condition and will
      not on its own make a process stop.
      The Conclusion had called it the sharpest result in its part,
      and two results in that part depend on it.
\item The fibration $p\colon\Hyp\to\catGSLT$ --- the object the
      Prestige's first half is really about --- has not been
      rigorously exhibited.
      Open problem~(ii) of Section~\ref{ndo:sec:open} says so.
      What the part does exhibit is a tower
      (Chapter~\ref{sec:fibration}), and
      Remark~\ref{rmk:not-a-fibration} concedes that the tower is an
      indexed family with overdetermined variance rather than the
      fibration earlier drafts called it.
\item Conjecture~\ref{conj:conjugacy} and
      Conjecture~\ref{conj:reversal}, on energy--time conjugacy
      and on reversal symmetry carrying the generator, were
      demoted from claims to conjectures in this draft.
\item Conjecture~\ref{conj:dissipation-choice}, relating
      dissipation at a cut to the strength of the choice needed to
      schedule it, is one-directional at best.
\item Conjecture~\ref{suf:conj:selected}, that internalization is
      selected above a threshold in the variance of a learner's
      inflow, is the price of the book's account of suffering and it
      has not been checked.
      The machinery to check it exists --- Chapter~\ref{ch:weight}
      supplies propensities and a working simulator --- so this is a
      gap that could be closed by running something rather than by
      proving something, which makes it the cheapest item on this
      list and the most embarrassing to have left open.
\item Proposition~\ref{prop:sim-noup}, that no position in the
      lattice can faithfully simulate a position above it, is used in
      Chapter~\ref{ch:sim} to rule out one form of the simulation
      hypothesis.
      It depends on subject reduction for schedulers, which is
      Gap~2 of Chapter~\ref{sec:hyper-gaps} and is open.
\item Conjecture~\ref{conj:fd}, that a budget supplies finite
      differentiation, is what the Prestige's account of notation
      rests on, and its transitivity is unverified.
\item Question~\ref{q:gsos}, whether the import of a world's
      behaviour into parallel composition is a distributive law.
      If it is not, compositionality does not transfer, and a
      good deal of the Prestige's construction goes with it.
\item Condition~\ref{ox:cond:rpo} and Condition~\ref{ox:cond:cong}
      --- that a theory has the relative pushouts its labels are
      defined by, and that the resulting bisimulation is a
      congruence for a stated class of contexts.
      These are assumed by every chapter of the machinery part and
      verified for no theory in this book.
      They are being established elsewhere for the theories that
      matter here~\cite{BIHOL}, and if they fail for a theory, the
      derived transition system that chapter uses does not exist.
      This is the most widely relied-upon unproved thing in the
      book, and it was invisible until Chapter~\ref{ch:obsext}
      forced it into the open.
\item Theorem~\ref{ox:thm:main} is proved for the first-order
      fragment only.
      Remark~\ref{ox:rem:binding} says what a binding constructor
      would need and does not supply it.
      Every calculus this book cares about has binders, so the
      theorem's coverage of the cases of interest is a hope rather
      than a result.
\end{itemize}

\noindent
Chapter~\ref{sec:quantum} is a fifth case of a different kind: it
is a chapter that reaches an obstruction and then goes sideways,
and it is marked as unfinished in its opening remark rather than at
its end.
Two of its results are stated without their proofs --- the
conservativity theorem over the summation-free stochastic
$\pi$-fragment, and the degeneration of the quantum-jump
unravelling to the Gillespie algorithm --- and both are proved in
the source note rather than here.
That is a deliberate compression and not a gap, but a reader who
wants to check them will have to go and get them.

\section*{The interludes, and the lowercase i}

Four short pieces of fiction sit between the parts, set in a
different typeface and a narrower measure: \emph{The Permitted Say}
before the Pledge, \emph{The Foam-Born} before the Turn,
\emph{begotten not made} before the Prestige, and
\emph{Right-Sizing} at the end.
They were written collaboratively with Claude, an AI system built
by Anthropic, and a note at the back of the book explains the
arrangement and names the debts.

Two conventions, since you will meet them before that note.
Throughout the main text i use a lowercase \emph{i} for myself; it
is not a typographical accident and the Pledge explains why.
Inside the interludes the first person is capitalised, because
inside the interludes the first person is not me.

\section*{What this book does not do}

Briefly, and without hedging.

It does not show that the mind is computational.
That premise is assumed on the first page and it is doing all the
work; Chapter~\ref{ch:trick} is about how much work.

It does not build the encoder --- the thing that would connect a
computational learner to a world not made of computations.
That is where the difficulty of any actual system would be
concentrated, and Chapters~\ref{ch:trick}
and~\ref{ch:notation} are an extended account of why.

It does not derive quantum mechanics.
Chapter~\ref{sec:quantum} gets as far as two obstructions, names
them, and then finds one place a presentation can put coherence
that neither obstruction touches --- the equations it declines to
quotient.
That is a smaller thing than a derivation and it is the honest size
of what is there.

It does not supply a metric on behaviours, without which the space
read off a term structure stays combinatorial rather than
becoming somewhere you could do physics.
That is stated as the next problem and it is not solved here.

And it is not a theory of phenomenal consciousness, whatever the
chapter titles may suggest.
It makes a claim about suffering --- a learner's own reading of its
distance from a level it holds as a goal --- and a claim about
consciousness as a position in a lattice of choice strength, and
neither claim answers the question of why there is something it is like
to be anything.
The first of those two words is chosen deliberately and it is heavier
than the structure it names; Chapter~\ref{ch:suffer} says so in its own
closing section, and i would rather be held to the structure.
The Prestige is candid about which of the book's questions have
been answered and which have only been made tractable.
Rather more of them are in the second category.

\bigskip
\noindent
That is the map.
The Turn begins on the other side of a story about a woman being
interviewed about something she saw, and it begins in earnest a few
pages after that.

%% restore the book's numbering for the numbered chapters that follow
\renewcommand{\thefigure}{\thechapter.\arabic{figure}}\setcounter{figure}{0}
